\chapter{Introduction} \label{chap:introduction}

\section{History of the problem}

It's 1956. The genius mathematician John von Neumann is hospitalized at the Walter Reed Army Medical Hospital in Washington, USA. One year earlier, a mass was found near his collarbone, which turned out to be cancer, probably caused by exposure to radiation at the Los Alamos National Laboratory during World War II \cite{vn_biography, vn_book}. After his condition improved thanks to a successful treatment, his colleague and friend Kurt Gödel sent a letter to express solidarity for his diagnosis, hoping for a full recovery. To keep von Neumann busy during his hospital stay, Gödel proposed his friend a seminal problem of fundamental interest both for logic and computer science.

\begin{quote}
    \say{One can obviously easily construct a Turing machine, which for every formula $F$ in first order predicate logic and every natural number $n$, allows one to decide if there is a proof of $F$ of length $n$ (length = number of symbols). Let $\psi(F,n)$ be the number of steps the machine requires for this and let $\varphi(n) = \max_F \psi(F,n)$. The question is how fast $\varphi$(n) grows for an optimal machine. One can show that $\varphi(n) \ge kn$.
    \newline
    If there really were a machine with $\varphi(n) \sim kn$ (or even $\sim kn^2$), this would have consequences of the greatest importance. Namely, it would obviously mean that in spite of the undecidability of the Entscheidungsproblem, the mental work of a mathematician concerning Yes-or-No questions could be completely replaced by a machine. After all, one would simply have to choose the natural number $n$ so large that when the machine does not deliver a result, it makes no sense to think more about the problem.}

    \hfill Gödel to von Neumann, Princeton, 20 March 1956 \cite{godel_letter, sipser_pnp}
\end{quote}

In a less formal way, Gödel's question asks if there is an efficient algorithm capable of deciding if a given mathematical statement has a proof of some given length. In modern times, the above problem became known as the \textit{Proof Size Problem} (PSP) \cite{psp_res, psp_res_2, pap}. Sadly, no reply from Von Neumann is known. During his long hospital stay, von Neumann's mental health rapidly declined due to his extreme thanatophobia, the fear of impending death \cite{vn_biography}. Almost a year later after being hospitalized, von Neumann died \cite{vn_book, vn_biography}.

When restricted from first-order predicate logic to propositional logic, Gödel's question becomes an historic precursor of the \textit{$\mathsf{P}$ vs.\,$\mathsf{NP}$} question, which asks whether problems whose solutions are easy to double-check, i.e. problems that lie in the class $\mathsf{NP}$, are also easy to solve from scratch, i.e. if they also lie in the class $\mathsf{P}$ \cite{cook_sat, levin_fsat}. The most interesting problems in $\mathsf{NP}$ are the so-called $\mathsf{NP}$-complete problems, that being problems in $\mathsf{NP}$ whose efficient solution would imply an efficient solution for all $\mathsf{NP}$ problems, i.e. $\mathsf{P} = \mathsf{NP}$. The current state of the art conjectures that $\mathsf{P} \neq \mathsf{NP}$, meaning that these so-called complete problems cannot be efficiently solved in their general case (see \Cref{sec:complexity} for more detailed information).

If this were the case, no algorithm could be substantially better than something based on a \textit{brute force} approach, the idea of testing every candidate solution in the search space defined by the problem itself. For instance, if a problem asks to determine if there is a sub-structure with a specific property, the brute force approach tests every possible sub-structure. In his letter, Gödel himself states that he sees no better approach than that of reducing the search space through clever techniques. Over the years, researchers began to explore the problem for \textit{propositional proof systems} \cite{cook_proof_compx, kraj_book, galesi_toran}, i.e. logical systems that use restricted rules to find proofs of propositional statements, confirming Gödel's hypothesis \cite{psp_res, psp_res_2}.

Although $\mathsf{P} \neq \mathsf{NP}$ would rule out efficient proof-search algorithms, the practical demand for Automated Theorem Proving (ATP) \cite{handbook_sat} kept growing as society advanced. This branch of computer science is applied in all areas where absolute reliability is required (hardware design, software verification, modern mathematics, \dots). To address this issue, computer scientists and logicians started investigating a closely related problem, referred to as automatability \cite{fis_int_weak_aut,res_wphard}. \textit{Automatability} asks if for a fixed proof system one can find an algorithm that returns a proof of a given statement with an at most polynomial blow-up with respect to the size of the shortest proof.

A central discovery of modern proof complexity is that natural propositional proof systems are not automatable under some widely believed complexity assumption, such as standard cryptographic assumptions \cite{kraj_pud_crypto, fis_int_weak_aut,ac0_crypto} or even $\mathsf{P} \neq \mathsf{NP}$ \cite{res_nphard,alg_nphard,cp_nphard,resk_nphard,ac0f_nphard}. Previously, similar results were known under stronger hypotheses \cite{res_wphard, res_hard_proofs,pc_wphard}. These results highlight a certain paradox in the world of proof search. One would expect that a more expressive proof system would be capable of achieving remarkably short proofs, making it easier to generate proofs. However, the very same increase in expressive power implies that the search space of those proofs may be so unstructured that no algorithm can find them faster than brute force. In other words, it seems that if an algorithm fails to discover a short proof that is guaranteed to exist then the fault lies not in the algorithm, but in the proof space itself.

These findings culminate into the brand new concept of analyzability \cite{pap}. While automatability evaluates whether an algorithm can efficiently output a proof bounded by the size of the shortest existing proof in a fixed system, \textit{analyzability} inspects the structural complexity of its search space. In particular, we are interested in understanding if proofs of lower bounds can leak information about the formula. For instance, we may be able to learn some information about the satisfiability of a formula by investigating what allows a proof system to achieve short proofs of the formula's lower bounds.

\section{Thesis structure and contributions (TODO)}

\cleardoublepage